\subsubsection{Test di Integrazione backend}
\label{subsec:test_integrazione_backend}
\begin{tabularx}{\textwidth}{|c|X|c|}
    \hline
    \textbf{ID} & \textbf{Descrizione} & \textbf{Stato} \\
    \hline
    \label{ti:ti9} T.I.9 & verificare che quando viene eseguita la GET dei dataset accadano le seguenti condizioni:
    \begin{itemize}
        \item nel caso in cui viene chiamata una get nell'endpoint \texttt{/datasets} verificare che il SqlAlchemyDatasetAdapter esegua il metodo corretto e che il backend risponda con un DatasetListDTO che contiene la lista di tutti i dataset salvati.
    \end{itemize} & Superato \\
    \hline 
    \label{ti:ti10} T.I.10 & verificare che quando viene eseguita la POST dei dataset accadano le seguenti condizioni:
    \begin{itemize}
        \item nel caso in cui viene chiamata una post nell'endpoint \texttt{/datasets} verificare che il SqlAlchemyDatasetAdapter esegua il metodo corretto e che il backend risponda con un DatasetDTO che contiene i metadati del dataset appena creato.
    \end{itemize} & Superato \\
    \hline
    \label{ti:ti11} T.I.11 & verificare che quando viene eseguita la PUT dei dataset accadano le seguenti condizioni:
    \begin{itemize}
        \item nel caso in cui viene chiamata una put nell'endpoint \texttt{/datasets/<dataset\_id>} verificare che il SqlAlchemyDatasetAdapter esegua il metodo corretto e che il backend risponda con un DatasetDTO che contiene i metadati del dataset appena aggiornato.
    \end{itemize} & Superato \\
    \hline  
    \label{ti:ti12} T.I.12 & verificare che quando viene eseguita la DELETE dei dataset accadano le seguenti condizioni:
    \begin{itemize}
        \item nel caso in cui viene chiamata una delete nell'endpoint \texttt{/datasets/<dataset\_id>} verificare che il SqlAlchemyDatasetAdapter esegua il metodo corretto e che il backend risponda con uno stato in base al risultato dell'operazione di cancellazione del dataset selezionato dal database.
    \end{itemize} & Superato \\
    \hline  
\end{tabularx}
